\documentclass[twoside]{article}
\usepackage{package}
\usepackage[T1]{fontenc}
\usepackage{subfiles}
% Silenciar la advertencia de versalita + negrita inexistente
\WarningFilter{latex}{Some font shapes were not available}
\WarningFilter{latex}{Font shape `T1/lmr/bx/sc'}

% Silenciar la advertencia de caracteres perdidos en el nullfont (para 01_introduccion.tex)
\WarningFilter{latex}{Missing character: There is no}
%Nuevos comandos
\renewcommand{\thefootnote}{\arabic{footnote}}
\raggedbottom



\begin{document}
\begin{center}
    \thispagestyle{empty}
    \large{\textsc{Características de sistemas lineales y Errores y ayuda en Matlab y Simulink}}\\[0.5cm] 
    \large{Fabrizio Benavidez\footnote[1]{\texttt{fabrizio.benavidez@mi.unc.edu.ar}}}\\[0.5cm]

    \large{\footnotemark[1]\textit{Facultad de Ciencias Exactas, Físicas y Naturales, Ingeniería Aeronáutica}}
    \\[0.5cm]

    \begin{minipage}{0.7\textwidth}
    \hrule
    \vspace{0.2cm}
        \normalsize{\textsc{Resumen}. Este trabajo abarca el análisis y la simulación de sistemas dinámicos lineales mediante MATLAB y Simulink. A través de un enfoque práctico, se explora desde la caracterización matemática de la respuesta de un sistema hasta la resolución de errores habituales de implementación y la presentación de resultados. Se realizó ejercicios del comportamiento dinámico mediante el análisis de ceros, polos, ganancia y sobrepaso. Se utilizan las funciones \texttt{zpk}, \texttt{stepinfo} e \texttt{impulse} para evaluar la respuesta temporal ante entradas tipo escalón. También se realizó la configuración y personalización para graficar figuras utilizando \texttt{plot}. Por último, se corrigió un código de MATLAB y un sistema generador en Simulink, identificando los errores y su posterior corrección.}\\[0.2cm]
        
        \normalsize{\textit{Palabras claves}: Sistemas Lineales, Ceros, Polos, Ganancia, Errores, Matlab, Simulink\ }
    \vspace{0.2cm}
    \hrule
    \end{minipage}
    
    \vfil

    \includegraphics[width = 0.4\textwidth]{imagenes/Matlab_Logo.png}\\[0.5cm]

    \vfil

    \today\\[0.2cm]
    \large{Semana 3}
\end{center}




\newpage

\section{Guía 05} 
\subfile{secciones/guia05}


\newpage

\section{Guía 06} 
\subfile{secciones/guia06}












\newpage
\thispagestyle{empty}
\bibliography{biblioteca.bib}
\bibliographystyle{plain}


\end{document}